%------------------------------------------------------------------------------%
\begin{question}
  Descreva o domínio da função
  \(
    f(x, y) = \ln\left(x^2 + y^2 - 4\right)
  \)
  \begin{enumerate}
    \item Qual a condição que os pontos do domínio deve satisfazer?
    \item Que região é essa? Represente-a graficamente.
    \item Qual o interior da região?
    \item Qual a fronteira da região?
    \item A região é aberta ou fechada?
    \item A região é limitada?
  \end{enumerate}
\end{question}

%------------------------------------------------------------------------------%
\begin{resolution}{Solução}
  \onslide<+->{Condição}
  \begin{align*}
    \onslide<+->{x^2 + y^2 - 4 & > 0 \\}
    \onslide<+->{x^2 + y^2 & > 4}
  \end{align*}

  \onslide<+->{Pontos exteriores ao circulo centrado na origem e raio 2}

  \onslide<+->{
  Domínio
  \[
    D_f = \left\{ (x, y)\in\R^2 \st x^2 + y^2 > 4\right\}
  \]}
\end{resolution}

%------------------------------------------------------------------------------%
\begin{resolution}{Representação Gráfica}
  \centering
  \includegraphics{ex_dominio_ln_x2_y2_4}
\end{resolution}

%------------------------------------------------------------------------------%
\begin{resolution}{Solução}
  \begin{enumerate}
    \item \onslide<+->{Interior da região}
    \onslide<+->{
    \[
      \text{Interior } = \left\{ (x, y)\in\R^2 \st x^2 + y^2 > 4\right\}
    \]}
    \vspace{-\baselineskip}
    \item \onslide<+->{Fronteira da região}
    \onslide<+->{
    \[
      \text{Fronteira } = \left\{ (x, y)\in\R^2 \st x^2 + y^2 = 4\right\}
    \]}
    \vspace{-\baselineskip}
    \item \onslide<+->A região é aberta ou fechada? \\
    \onslide<+->{Aberta,}
    \onslide<+->{todos os pontos são pontos interiores}
    \item \onslide<+->{A região é limitada?} \\
    \onslide<+->{Ilimitada,}
    \onslide<+->{pois dado qualquer disco sempre teremos um ponto fora dele}
  \end{enumerate}
\end{resolution}

%------------------------------------------------------------------------------%